\paragraph{VLA baseline implementation and fine-tuning.}
\label{par:vla_reproducibility}
Ten VLA baselines are adapted to Touch, Pick, Move and Put with UR5e--Robotiq
2F-85 in PyBullet. Tables~\ref{tab:vla_adaptation}--\ref{tab:vla_optimizers} specify
trainable components, update budgets and optimization.
Exact weights, split IDs, calibration and recipes are fixed in
\texttt{vla21-lock-v1.json}.

\noindent\textit{Data.}
Hashed scene-group splits reserve 2,107 original training and 477 validation
trajectories. $D_2$ adds 38 successful corrections; $D_1$ adds 40 pose-recovery
trajectories to $D_2$, and $D_3$ adds 152 near-miss corrections to $D_1$.
Deviated states receive expert action labels; perturbation commands are excluded.
Original demonstrations are retained. Tasks/targets are
balanced; FLOWER main/$\pi_0$/VLA-JEPA use original/corrective/pose source weights
0.50/0.15/0.35, renormalized over available sources. $D_3$ near-miss replacement
has probability 0.35; within it, the first eight frames receive probability 0.60.
Images/state are aligned before each action; statistics use training data only.
Final-budget checkpoints
are evaluated without test-based selection.

\begin{table}[!ht]
\centering\footnotesize
\caption{VLA adaptation protocol. All stages use the same 21-D shared state.
Batch is the final-stage effective batch; updates accumulate across stages.
$H/K$: predicted/executed chunk length; commands run at 10\,Hz.}
\label{tab:vla_adaptation}
\begin{tabularx}{\linewidth}{@{}l c Y c r c c@{}}
\toprule
Method & Data & Trainable components & Batch & Updates & Output & $H/K$\\
\midrule
FLOWER~\cite{reuss2025flower} & $D_1$ & VLM LoRA ($r=8$), flow/action core, state projection & 16 & 4,224$^a$ & $\Delta7$ & $10/10$\\
X-VLA~\cite{zheng2026xvla} & $D_2\!\to\!D_1$ & All-linear LoRA ($r=8$), soft prompts, action encoder/decoder & 4 & 14,500$^b$ & Abs20$\to\Delta7$ & $30/30$\\
$\pi_0$~\cite{black2025pi0} & $D_1$ & Language/action LoRA ($r=16/32$), state/action/time projections & 4 & 3,000 & $\Delta7$ & $5/4$\\
VLA-Adapter~\cite{wang2026vla} & $D_2$ & All-linear LoRA ($r=64$), action queries, $L_1$ head, proprio projector & 4 & 6,000 & $\Delta7$ & $8/8$\\
SmolVLA~\cite{shukor2025smolvla} & $D_2$ & Language/connector, action expert and state/action projections; vision frozen; no LoRA & 8 & 22,500$^c$ & $\Delta7$ & $5/5$\\
VLA-JEPA~\cite{sun2026vla} & $D_1$ & Action-only: DiT/action network and state encoder; Qwen frozen & 4 & 2,500 & $\Delta7$ & $7/7$\\
Dita~\cite{hou2025dita} & $D_1$ & Full DINO/Q-Former/diffusion policy and sensor projection; CLIP text encoder frozen; no LoRA & 32 & 10,000$^d$ & $\Delta7$ & $10/1$\\
VLANeXt-S~\cite{wu2026vlanext,wu2026vlanextfamily} & $D_1\!\to\!D_3$ & Full native policy: VLM, flow/action core and proprio projector; no LoRA & 32$^e$ & 9,728$^e$ & $\Delta7$ & $8/1$\\
FALCON-ESM~\cite{zhang2026falcon} & $D_3$ & VLM LoRA ($r=8$), spatial injector, action/gripper heads and state projection; ESM frozen & 16 & 2,050$^f$ & $\Delta7$ & $10/1$\\
Seer~\cite{tian2025seer} & $D_3$ & Native transformer, resampler, action/image-prediction heads and state projections; CLIP/MAE frozen; no LoRA & 16 & 4,098$^g$ & $\Delta7^h$ & $3/1$\\
\bottomrule
\end{tabularx}
\vspace{2pt}
\begin{minipage}{\linewidth}\scriptsize
$D_1$: 2,185 trajectories /119,495 timesteps;
$D_2$: 2,145 /117,736 ($D_1$ without 40 pose-recovery trajectories);
$D_3$: 2,337 /125,527 ($D_1$ +152 near-miss corrections).
All exclude 477 validation trajectories; synchronized camera streams share timestep counts.
$^a$4,096 main +128 rehearsal; two preceding qualification updates excluded.
$^b$7,000 on $D_2$ +7,500 on $D_1$.
$^c$20,000 preceding-stage +2,500 final-stage updates; final batch 8.
$^d$Checkpoint at 10,000 updates.
$^e$8,192 on $D_1$ +1,024 recovery +512 rehearsal on $D_3$; initial batch 16, final batch 32.
$^f$2 qualification +2,048 adaptation.
$^g$2 qualification +4,096 adaptation.
$^h$Seer uses previous-requested-pose-anchored local increments, converted to $\Delta7$.
\end{minipage}
\end{table}


\noindent\textit{Observations.}
Each policy receives the instruction, external/wrist RGB and an ordered 21-D
robot/interaction vector: joint angles (6), measured gripper closure (1),
pad-centre position (3), preceding gripper command (1), tool rotation (6), jaw
aperture (1), left/right pad-contact indicators (2) and total pad normal force (1).
Positions use world-frame metres, joints radians and force newtons; tool rotation
stacks the first two tool-to-world matrix columns. Closure is
in $[0,1]$ and the preceding command in $[-1,1]$, positive for closing.
At each query, each contact bit is one if that pad's summed external-contact
normal force exceeds 0.1\,N; the force channel sums both pads without object
identity filtering. Continuous channels use training mean/std with a
$10^{-6}$ standard-deviation floor; the two contact bits remain unchanged.

FLOWER, SmolVLA and VLA-JEPA project the 21-D vector directly. X-VLA prepends
its native 20-D state (10 active-arm pose/closure channels and 10 zeros), giving
41 dimensions. VLA-Adapter prepends end-effector XYZ, rotation vector and
half-apertures $(d/2,-d/2)$ ($d$: jaw aperture), giving 29 dimensions.
Dita/VLANeXt-S prepend XYZ, rotation vector and closure (28-D total).
Seer projects a separate 7-D XYZ/Euler/open-state
prefix and the 21-D vector. Native rotation/preprocessing conventions are
retained. $\pi_0$ pads the normalized state with eleven zeros and serializes
the same raw 21 fields to four decimals. Native prefix weights are preserved;
new conditioning columns, or the output layer of a new residual branch,
are zero-initialized.

Only two RGB views are supplied. FALCON-ESM also receives metric front depth and
camera calibration, transformed with image crop/resize. RGB inputs are
$384\times384$, processed to $224\times224$ except SmolVLA ($512\times512$)
and VLANeXt-S ($256\times256$). Dita uses two image/state timestamps, Seer seven,
and VLANeXt-S eight state timestamps with current images; others use one.
History uses 10\,Hz samples including the current observation; reset padding
repeats the initial observation.

\noindent\textit{Initialization and training.}
Released checkpoint IDs and hashes are fixed in the lock. Notable choices are
FLOWER's 12-block variant, VLA-Adapter LIBERO-Spatial-Pro and VLANeXt-S's
Qwen3.5-0.8B initialization rather than pretrained robot action modules.
Unlisted components are frozen. LoRA uses $\alpha/r=2$ for
FLOWER, X-VLA, VLA-Adapter and FALCON-ESM.
X-VLA trains LoRA/prompts/action modules jointly instead of its published
initial prompt/head-only stage; VLA-JEPA uses
action-only training with frozen Qwen, no world-model loss and standardized
Euler-delta targets rather than min--max-scaled axis--angle targets.

Flow policies use action flow matching; VLA-Adapter uses $L_1$;
Dita uses masked diffusion-noise mean-squared error (MSE); FALCON-ESM uses masked arm Huber plus
0.01 gripper binary cross-entropy (BCE) with logits. Seer uses arm smooth-$L_1$
plus 0.01 gripper BCE and 0.1 image MSE at $t+3$; future images are training
targets only. VLANeXt-S adds discrete-cosine-transform regularization (0.1).
Dita uses brightness jitter, VLANeXt-S crop/colour jitter and Seer temporally
consistent shifts; ranges are fixed in the lock. Inference uses four
sampling steps for FLOWER/VLA-JEPA and ten for X-VLA/$\pi_0$/SmolVLA/VLANeXt-S;
Dita uses 100 training noise levels and ten inference steps.
$\pi_0$ uses JAX/OpenPI; VLANeXt-S uses AdamW8bit with mixed BF16/FP32 parameters;
other trainable master weights/moments are FP32.

\begin{table}[!ht]
\centering\footnotesize
\caption{AdamW settings ($\beta_1=0.9$, $\epsilon=10^{-8}$).
$B_\mu\times A$: per-step batch times gradient accumulation; LR: learning rate; WD: weight decay;
clip: global gradient norm. Schedules count optimizer updates unless noted.}
\label{tab:vla_optimizers}
\begin{tabularx}{\linewidth}{@{}l c l c l c Y@{}}
\toprule
Method/stage & $B_\mu\times A$ & Peak LR & $\beta_2$ & WD & Clip & Warmup / schedule\\
\midrule
FLOWER main & $1\times16$ & $10^{-5}/2\!\times\!10^{-5}$ & .95 & .05 & 1 & 128; cosine to half peak at 4,096\\
FLOWER rehearsal & $1\times16$ & $2\!\times\!10^{-6}/4\!\times\!10^{-6}$ & .95 & .05 & 1 & Constant; 128 updates, fresh optimizer\\
X-VLA & $1\times4$ & $10^{-4}$ & .99 & 0 & 10 & 200; cosine to 0.1 peak; final horizon 14,500\\
$\pi_0$ & $1\times4$ & $2.5\!\times\!10^{-5}$ & .95 & $10^{-10}$ & 1 & 100; cosine to $2.5\!\times\!10^{-6}$ at 3,000\\
VLA-Adapter & $1\times4$ & $10^{-4}$ & .99 & 0 & 10 & 200; cosine to $10^{-5}$ at 6,000\\
SmolVLA final & $2\times4$ & $2.5\!\times\!10^{-5}$ & .95 & $10^{-10}$ & 10 & 600; cosine over 10,000 microsteps to $6.25\!\times\!10^{-7}$\\
VLA-JEPA & $1\times4$ & $10^{-4}$ & .95 & $10^{-10}$ & 10 & 100; cosine to $10^{-5}$ at 2,500\\
Dita & $1\times32$ & $10^{-4}$ & .95 & .05 & 1 & 100; cosine over 10,000 to 0.01 peak; DINO LR $\times0.1$\\
VLANeXt-S main/recovery & $1\times16$ & $10^{-4}$ & .999 & .1 & 1 & 1,000; cosine horizon 20,000, stopped at 9,216\\
VLANeXt-S rehearsal & $1\times32$ & $2\!\times\!10^{-5}$ & .999 & .1 & 1 & 16; 512-update extension with $10^{-6}$ floor\\
FALCON-ESM & $1\times16$ & $5\!\times\!10^{-5}$ & .999 & 0 & 1 & 64; cosine horizon 4,096 to 0.1 peak; selected at 2,048 adaptation updates\\
Seer & $1\times16$ & $10^{-4}$ & .999 & $10^{-4}$ & 1 & 64; cosine horizon 4,096 to 0.1 peak\\
\bottomrule
\end{tabularx}
\vspace{2pt}
\begin{minipage}{\linewidth}\scriptsize
FLOWER LR pairs denote VLM/non-VLM groups. Dita excludes biases/norms from WD.
VLANeXt-S uses AdamW8bit. FALCON-ESM/Seer have two preceding qualification
updates at batch 4, LR $10^{-4}$, WD $10^{-4}$.
SmolVLA's stages use 10,000/5,000/5,000/2,500 updates at batches 32/32/16/8;
the first-stage peak/end LR is $5\!\times\!10^{-5}/1.25\!\times\!10^{-6}$,
then $2.5\!\times\!10^{-5}/6.25\!\times\!10^{-7}$.
Its final 600/10,000 microsteps equal 150/2,500 optimizer updates.
Earlier warmups are 1,000/250/250 updates.
\end{minipage}
\end{table}


\noindent\textit{Execution.}
The shared $\Delta7$ action is world-frame XYZ translation (m), local XYZ Euler
increments (rad) and gripper command ($-1$ open, $+1$ close). Arm-action scaling
uses training 1st/99th percentiles for FLOWER/VLA-Adapter/Dita/VLANeXt-S,
mean/std for $\pi_0$/SmolVLA/VLA-JEPA, and percentiles mapped to $[-0.65,0.65]$
for FALCON-ESM. Native gripper conventions are converted to this interface.
X-VLA predicts absolute 20-D pose/closure; the active arm is converted against
the measured end-effector pose at each executed step. Seer's local increments
are anchored to the preceding requested pose, not the current measured pose.
The command rate is 10\,Hz with 240\,Hz physics (24 substeps/command).
$H$ denotes predicted chunk length and $K$ the number executed before querying
again; thus the query rate is $10/K$\,Hz. Remaining predictions are discarded;
no temporal ensembling is used. The shared inverse-kinematics servo clips
translation/rotation increments to $\pm0.05$\,m/$\pm0.25$\,rad and retains the
common safety checks and contact-qualified grasp constraint; executor sources
are hash-bound in the lock.
